\chapter{Related Work}
\label{ch:related}

This chapter reviews earlier work from the view of psychiatric decision support. Many psychiatric LLM studies are evaluated mainly by the final diagnosis. Some systems also provide explanations, criterion records, or intermediate decisions, but these outputs are not always evaluated as separate parts of the diagnostic process.

We review five topics: psychiatric LLM diagnosis; decision support, auditability, and criterion-level records; primary-diagnosis selection; retrieval with similar cases; and medical multi-agent systems. The goal is to identify what information each system provides, how that information is evaluated, and which parts of the diagnostic process remain difficult to study.

\section{Psychiatric LLM Diagnosis}

Psychiatric LLM studies cover many different tasks. Some test psychiatric knowledge or licensing questions. Others use short clinical cases, focus on one disorder, detect depression from text, diagnose from a fixed record, or support an interactive consultation~\citep{qiu2023largeai,li2024taiwan,raballo2025,sarma2025,li2025addreview,mcduff2025amie}. These tasks differ in the evidence given to the model, the number of possible diagnoses, the required output, and the label used for evaluation.

Screening is different from choosing among several similar diagnoses. A screening system may only need to find a broad depression or anxiety signal. A differential-diagnosis system must compare disorders that share symptoms and still choose one primary diagnosis. This choice may depend on symptom duration, which symptoms started first, illness course, other possible causes, and comorbidity. LingxiDiagBench reports lower performance as the task moves from broad depression--anxiety classification to comorbidity recognition and twelve-category diagnosis~\citep{xu2026lingxidiagbench}. LingxiDiagBench and MDD-5k provide synthetic Chinese benchmarks for multi-category psychiatric diagnosis~\citep{xu2026lingxidiagbench,yin2024mdd5k}.

Another key difference is whether the system can ask more questions. MIND, WiseMind/ProAI, MAGI, DSM5AgentFlow, and conversational diagnostic AI study interactive workflows that can ask questions during the assessment~\citep{li2026mind,wu2026wisemind,bi2025magi,ozgun2025dsm5agentflow,tu2025conversational}. MentalSeek-Dx and MoodAngels instead work with information that is already present in a record or transcript~\citep{sun2026mentalseek,xiao2025moodangels}. An interactive system can collect new evidence. A fixed-input system must work with the same possibly incomplete record.

Taken together, earlier work shows that LLMs can produce psychiatric diagnoses in several settings. However, many of the studies reviewed in this chapter are evaluated mainly by the final diagnosis. Some also provide explanations or intermediate records, but they do not separately measure whether a gold label was proposed as a candidate, supported by the criterion record, or selected as primary. This makes it difficult to identify where a disagreement appeared and what information would be useful for clinical review.

\section{Decision Support, Auditability, and Criterion-Level Records}
\label{sec:auditability-definition}

For clinical decision support, a system should provide more than one final diagnosis. Useful outputs may include possible diagnoses, supporting and conflicting evidence, criterion states, and information that is still missing. These outputs can help a reviewer understand what the system considered and what should be checked next.

Diagnostic-chain and symptom-based methods can show intermediate diagnoses, confidence values, or clinically motivated features~\citep{chen2025cod,lee2026interpretable,lyu2026depllm}. Psychiatric systems may also connect outputs to guidelines, structured interviews, symptom-to-criterion links, and supporting or conflicting evidence~\citep{li2026mind,bi2025magi,ozgun2025dsm5agentflow}. Constraint-logic methods make the rule layer easier to inspect by turning diagnostic criteria into explicit programs~\citep{kim2025clp}.

Visible output does not guarantee a correct or faithful explanation. A model may write a clear explanation that does not match the factors that produced its answer~\citep{turpin2023unfaithful}. A saved evidence note may be unsupported, and a visible rule may still contain a weak assumption. Explainability, interpretability, faithfulness, and auditability are therefore different properties in healthcare~\citep{rudin2019stop,huang2024xiai}.

In this thesis, auditability means that observable outputs are saved and can be checked after the system finishes. It does not mean that the system reveals its hidden reasoning or that the saved records are clinically correct. Expert review of model-generated diagnostic rules and evidence remains necessary~\citep{kim2025clp}.

The remaining question is whether these records are evaluated together on the same cases. A useful stage-wise evaluation should show whether a diagnosis was proposed, whether the recorded criteria kept it compatible with the transcript, and whether it was finally selected as primary.

\section{Primary Diagnosis Selection}

Criterion checking and primary-diagnosis selection answer different questions. Criterion checking asks whether one diagnosis is compatible with the available transcript under a defined rule. Primary-diagnosis selection asks which diagnosis best explains the case when several diagnoses remain possible.

This difference is important in psychiatry because several disorders may share low mood, worry, poor sleep, fatigue, poor concentration, and reduced daily function. More than one diagnosis may therefore pass its own criterion rule. A compatible diagnosis may be primary, comorbid, secondary, or only a differential alternative.

Structured medical knowledge can help collect or organize evidence. MedKGI combines a medical knowledge graph, guided questions, and a structured dialogue state to seek useful information during diagnosis~\citep{wang2025medkgi}. DKEC builds knowledge graphs from outside medical sources and uses their links for multilabel diagnosis prediction~\citep{ge2024dkec}. These studies show how structured information can support evidence collection or label prediction. Their tested goals are different from choosing one primary diagnosis among several criterion-compatible psychiatric candidates in a fixed transcript.

A primary selector may need to compare candidates directly. Useful information may include symptom timing, duration, illness course, past episodes, other possible causes, which syndrome explains the case best, and whether more than one disorder is present.

Among the studies reviewed in this chapter, none jointly reports and evaluates a ranked candidate list, criterion states, criterion-compatible diagnoses, and the final primary choice on the same cases. This makes it difficult to distinguish a missing candidate from a criterion disagreement or a primary-selection disagreement.

\section{Retrieval with Similar Cases}

Retrieval can add either knowledge sources or example cases. Knowledge retrieval adds documents, guidelines, or other factual sources~\citep{lewis2020rag}. Example retrieval places similar labeled cases in the prompt for in-context learning~\citep{liu2022incontext}. A retrieved example may guide the model, but it is not evidence about the current patient.

Example choice can affect in-context learning. Similar examples may show useful wording, expected output formats, and common label patterns. However, repeated frequent classes may strengthen corpus shortcuts, while weakly related examples may distract the model~\citep{liu2022incontext}.

A label-aware retrieval policy may show a wider range of diagnoses, but it also changes the number, order, similarity, total length, and mix of examples. Retrieval policies should therefore be compared within the same architecture rather than assumed to provide an automatic benefit.

\section{Medical Multi-Agent Systems}

A medical multi-agent system uses several model calls and gives them different roles. One call may propose a diagnosis, another may extract evidence, another may criticize an answer, and another may make the final decision. Other systems use consultation, debate, voting, or consensus. MedAgents uses several role-based medical experts, while MDAgents changes the collaboration plan according to estimated task difficulty~\citep{tang2024medagents,kim2024mdagents}. Multi-agent debate is a related method in which agents exchange and revise competing answers~\citep{du2024debate}.

These designs may help because a narrow role can focus on one task, a second model call may find missing evidence, and a structured workflow can save intermediate outputs for review. These are possible benefits, not guarantees. More agents or more discussion rounds do not automatically lead to better diagnostic accuracy.

Multi-agent systems can also repeat the same mistake. Agents that use the same model family, similar prompts, and the same evidence may fail in similar ways. A wrong candidate or an unsupported evidence statement can pass into later steps. Free-form discussion may also change several parts of the workflow at once.

This makes the source of a performance change hard to identify. If the candidate generator, evidence record, discussion process, and final rule all change together, a higher or lower final score is a result of the complete system. A comparison with a Single LLM may also change the prompt, context, number of model calls, compute, and output contract.

The value of a multi-agent design is therefore not limited to higher final accuracy. Separate roles can also make candidate generation, criterion checking, and final selection visible as different outputs. This value depends on whether those outputs are saved and evaluated, not simply on the number of agents or discussion rounds.

\section{Research Gap and Research Questions}

The main gap is not the absence of another final-label predictor. The main gap is the lack of a stage-wise decision-support evaluation that connects three questions on the same cases:

\begin{enumerate}
    \item Was at least one gold label proposed as a candidate?
    \item Did the criterion-level record place at least one gold label in the criterion-compatible set?
    \item Was at least one gold label selected as the primary diagnosis?
\end{enumerate}

Many reviewed systems provide only the final diagnosis, while others provide explanations or criterion-related records without separately evaluating these three decisions. As a result, one final score cannot show whether a disagreement appeared during candidate generation, criterion checking, or primary-diagnosis selection.

Table~\ref{tab:related-work-matrix} summarizes the outputs reported in representative studies. It describes the published evaluation of each study. It does not rank overall system quality or clinical ability.

\begin{table}[htbp]
\centering
\caption{Representative psychiatric diagnostic systems and the outputs reported for stage-wise review.}
\label{tab:related-work-matrix}
\small
\begin{adjustbox}{max width=\textwidth}
\begin{tabular}{|>{\raggedright\arraybackslash}p{2.5cm}|>{\raggedright\arraybackslash}p{2.7cm}|>{\raggedright\arraybackslash}p{3.1cm}|>{\raggedright\arraybackslash}p{4.3cm}|>{\centering\arraybackslash}p{2.0cm}|}
\hline
Study or system & Evidence setting & Candidate output & Criterion or evidence record & Stage-wise evaluation\\
\hline
MIND~\citep{li2026mind} & Interactive consultation & No ranked evaluation reported & Criterion-based rationales and process rewards & No\\
\hline
WiseMind/\allowbreak ProAI~\citep{wu2026wisemind} & Interactive consultation & No ranked evaluation reported & Structured DSM-guided action states & Partial\\
\hline
MAGI~\citep{bi2025magi} & Interactive MINI-guided interview & No ranked evaluation reported & Symptom and dialogue evidence linked to criteria & Partial\\
\hline
DSM5Agent\allowbreak Flow~\citep{ozgun2025dsm5agentflow} & Interactive questionnaire & No ranked evaluation reported & Utterance-level supporting and conflicting criteria & Partial\\
\hline
MentalSeek-Dx~\citep{sun2026mentalseek} & Fixed clinical record & Possible diagnoses generated; ranking not evaluated & Criteria-guided staged reasoning & Partial\\
\hline
Mood\allowbreak Angels~\citep{xiao2025moodangels} & Fixed records and assessment scales & No ranked evaluation reported & DSM-5 retrieval and symptom matching & Partial\\
\hline
Constraint-logic CDSS~\citep{kim2025clp} & Structured patient facts & None & Inspectable diagnostic rules & No\\
\hline
\end{tabular}
\end{adjustbox}

\vspace{0.4em}
\begin{minipage}{0.98\textwidth}
\footnotesize
\textit{Note.} ``Partial'' means that a study reports intermediate diagnostic or criterion-related outputs but does not separately evaluate candidate generation, criterion checking, and primary selection on the same cases. ``Not reported'' refers only to the published evaluation and does not show that the implementation lacks the capability.
\end{minipage}
\end{table}
\FloatBarrier

This thesis addresses the main gap through the stage-wise analysis in RQ2. The other research questions provide benchmark context or test supporting design choices. RQ1 compares complete systems on the same internal split. RQ3 studies similar-case retrieval within each architecture. RQ4 tests alternative primary-selection methods. RQ5 applies the same output views to a second synthetic source.

This thesis uses the following research questions:

\begin{description}[style=nextline,leftmargin=1.5cm,labelwidth=1.1cm,itemsep=0pt,parsep=0pt,topsep=0pt,partopsep=0pt]
\item[RQ1] How does HiED perform relative to conventional classification and a Single LLM on the same internal split and parent-label evaluation?
\item[RQ2] Where is benchmark disagreement recorded across candidate generation, criterion checking, compatibility analysis, and primary commitment?
\item[RQ3] How do similar-case retrieval strategies affect Single and HiED?
\item[RQ4] Do the tested primary-selection strategies reduce the recorded selection disagreement?
\item[RQ5] Does the stage-wise analysis remain useful under a second synthetic source?
\end{description}

RQ2 is the main stage-wise analysis. RQ1 provides benchmark context for the complete systems. RQ3 examines retrieval as a supporting configuration. RQ4 tests whether alternative selection methods reduce the recorded disagreement. RQ5 examines whether the same analysis can be applied to a second synthetic source.
